BIÊN BẢN PHỎNG VẤN
VỊ TRÍ: GIÁM ĐỐC SẢN PHẨM CAO CẤP (GIẢI PHÁP AI)
ỨNG VIÊN
Alex Morgan
NGƯỜI PHỎNG VẤN
Sarah Jenkins (VP Sản phẩm)
NGÀY \& GIỜ
Ngày 24 tháng 10 năm 2026
VÒNG
Giai đoạn 2 - Kỹ thuật \& Lãnh đạo
Bối cảnh tổng quan: Phỏng vấn vòng hai nhằm đánh giá mức độ phù hợp chiến lược, nền tảng kỹ thuật trong phát triển sản phẩm AI và năng lực lãnh đạo của ứng viên để dẫn dắt đội ngũ Nền tảng Đám mây Doanh nghiệp.

PHẦN 1: THÔNG TIN NỀN TẢNG \& TẦM NHÌN
Sarah Jenkins
Chào mừng Alex. Cảm ơn bạn đã dành thời gian hôm nay. Để bắt đầu, bạn có thể tóm tắt ngắn gọn kinh nghiệm dẫn dắt các sáng kiến AI cho doanh nghiệp và điều gì đã thu hút bạn đến với vị trí này tại NexusTech không?
Alex Morgan
Cảm ơn Sarah. Rất hân hạnh được ở đây. Trong sáu năm qua tại CloudScale Inc., tôi đã quản lý các sản phẩm B2B SaaS, đặc biệt tập trung vào bộ công cụ phân tích máy học. Tôi đã dẫn dắt quá trình chuyển đổi từ phân tích truyền thống sang các mô hình AI dự đoán, giúp tăng 35\% ARR. Tôi đã theo dõi sự mở rộng của NexusTech sang các giải pháp AI cho doanh nghiệp và rất ấn tượng với định hướng của công ty vào cơ sở hạ tầng có thể mở rộng. Tôi thấy sự phù hợp mạnh mẽ giữa kinh nghiệm tích hợp doanh nghiệp của tôi và lộ trình phát triển của công ty.
Sarah Jenkins
Đó là một thành tích rất ấn tượng. Bạn có thể đưa ra một ví dụ cụ thể về cách bạn ưu tiên các tính năng khi nợ kỹ thuật (technical debt) cạnh tranh trực tiếp với các yêu cầu từ khách hàng có nhu cầu cao không?
Alex Morgan
Tất nhiên rồi. Năm ngoái, dịch vụ lập chỉ mục cốt lõi của chúng tôi gặp phải nợ kỹ thuật từ hệ thống cũ, làm chậm tốc độ truy vấn gần 40\%. Cùng lúc đó, ba khách hàng doanh nghiệp lớn đang yêu cầu các tính năng báo cáo tùy chỉnh. Tôi đã áp dụng mô hình RICE có trọng số, phân bổ chính xác 20\% nguồn lực kỹ thuật (bandwidth) mỗi sprint cho việc tái cấu trúc kiến trúc, đồng thời đàm phán mốc thời gian bàn giao theo từng giai đoạn cho các yêu cầu của khách hàng. Trong vòng hai quý, độ trễ truy vấn đã giảm một nửa và chúng tôi đã giao 90\% tập tính năng được yêu cầu mà không gây kiệt sức cho đội ngũ kỹ sư.

PHẦN 2: CHIẾN LƯỢC SẢN PHẨM \& GIẢI QUYẾT VẤN ĐỀ KỸ THUẬT
Sarah Jenkins
Bạn tiếp cận như thế nào khi thiết kế các tính năng dựa trên AI mà đầu ra của mô hình mang tính xác suất thay vì định tính (nguyên tắc xác định)? Bạn quản lý kỳ vọng của người dùng và trải nghiệm người dùng (UX) xung quanh các lỗi như thế nào?
Alex Morgan
Thiết kế cho AI đòi hỏi sự chuyển dịch từ kết quả cố định sang khoảng tin cậy. Đầu tiên, chúng tôi thiết kế sự minh bạch trực tiếp vào giao diện người dùng (UI)—hiển thị điểm tin cậy hoặc giải thích lý do tại sao một đề xuất lại được đưa ra. Thứ hai, chúng tôi đảm bảo các phương án dự phòng có sự tham gia của con người (human-in-the-loop) một cách liền mạch. Nếu ngưỡng tin cậy giảm xuống dưới 80\%, UI sẽ gợi ý xem xét thủ công nhanh thay vì tự động đưa ra quyết định một cách âm thầm.
Đánh giá Tuyển dụng NexusTech • Ứng viên: Alex Morgan
Trang 1 / 2

PHẦN 2: CHIẾN LƯỢC \& GIẢI QUYẾT VẤN ĐỀ (TIẾP THEO)
Sarah Jenkins
Làm thế nào để bạn xác minh mức độ phù hợp giữa sản phẩm và thị trường (product-market fit) trước khi đầu tư nguồn lực kỹ thuật lớn vào một mô hình AI mới?
Alex Morgan
Tôi kiên quyết ủng hộ việc thử nghiệm theo phương pháp 'Wizard of Oz' hoặc các bản mẫu thử nghiệm độ trung thực thấp (low-fidelity prototypes). Trước khi xây dựng các mô hình tùy chỉnh hoàn chỉnh, chúng tôi thử nghiệm nhu cầu bằng cách sử dụng API của bên thứ ba hoặc các quy trình bán thủ công để quan sát mức độ tương tác của người dùng. Chúng tôi đo lường thành công bằng các chỉ số định tính chính như tốc độ hoàn thành quy trình công việc và độ hài lòng định tính, thay vì các chỉ số phù phím ban đầu.

PHẦN 3: NĂNG LỰC LÃNH ĐẠO ĐA CHỨC NĂNG
Sarah Jenkins
Hãy mô tả một tình huống mà bộ phận kỹ thuật không đồng ý với chiến lược sản phẩm của bạn. Bạn đã làm thế nào để thống nhất ý kiến giữa các đội ngũ?
Alex Morgan
Các Trưởng nhóm Kỹ thuật của chúng tôi muốn xây dựng lại toàn bộ quy trình (pipeline) bằng một framework mới, ước tính mất bốn tháng làm việc mà không tạo ra

giá trị hướng tới khách hàng. Tôi đã tổ chức một buổi workshop chung, nơi chúng tôi đối chiếu các điểm đau của nhà phát triển với dữ liệu khách hàng rời bỏ dịch vụ. Chúng tôi đã thống nhất về một chiến lược chuyển đổi từng bước, giải quyết các nút thắt cổ chai tồi tệ nhất trước tiên trong khi vẫn tiếp tục ra mắt các tính năng có tác động cao trong mỗi chu kỳ phát hành.

ĐÁNH GIÁ PHỎNG VẤN \& ĐIỂM SỐ ỨNG VIÊN

LĨNH VỰC NĂNG LỰC
TÍN HIỆU / BẰNG CHỨNG QUAN SÁT ĐƯỢC
ĐÁNH GIÁ

Nhạy bén Kỹ thuật
Hiểu biết sâu sắc về quy trình làm việc AI, tích hợp API và các sự đánh đổi trong hệ thống.
4.5 / 5.0

Chiến lược Sản phẩm
Thể hiện khả năng sắp xếp thứ tự ưu tiên có cấu trúc bằng cách sử dụng các khung (RICE, MVP khám phá).
4.8 / 5.0

Quản lý Bên liên quan
Phương pháp tiếp cận thực tế đối với sự đánh đổi về kỹ thuật và giao tiếp với khách hàng.
4.2 / 5.0

Lãnh đạo \& Văn hóa
Có tính hợp tác, dựa trên dữ liệu và thể hiện sự thấu cảm rõ ràng đối với các đội ngũ kỹ thuật.
4.6 / 5.0

Khuyến nghị của Người phỏng vấn: ĐỀ XUẤT CHUYỂN TIẾP MẠNH MẼ
Alex đã thể hiện chuyên môn sâu sắc về quản lý sản phẩm B2B cùng với khả năng nắm bắt xuất sắc về các hạn chế của sản phẩm AI. Khuyến nghị đưa ứng viên vào vòng phỏng vấn cuối cùng với Giám đốc Kỹ thuật (VP of Engineering) và Giám đốc Sản phẩm (CPO).

Đánh giá Tuyển dụng NexusTech • Ứng viên: Alex Morgan
Trang 2 / 2
